\documentclass[11pt,a4paper]{article}
\usepackage[utf8]{inputenc}
\usepackage[T1]{fontenc}
\usepackage[brazil]{babel}
\usepackage[a4paper,margin=2.3cm]{geometry}
\usepackage{enumitem}
\usepackage{hyperref}
\usepackage{xcolor}
\setlength{\parindent}{0pt}
\setlength{\parskip}{5pt}
\sloppy
\setlength{\emergencystretch}{3em}
\hypersetup{colorlinks=true,urlcolor=blue,linkcolor=black}
\newcommand{\cabecalho}[2]{\begin{center}{\small \textbf{Sprint 2 -- #1}}\\[1mm]\textit{Inventário de Ativos e Vulnerabilidades}\\Matheus Moreli Durão Silva -- 12621CBS217\\Cibersegurança -- Bacharelado -- Noturno\end{center}\hrule\vspace{4mm}}
\newcommand{\bloco}[1]{\par\vspace{2mm}\textbf{#1}\par}
\newcommand{\codigo}[1]{\texttt{#1}}

\begin{document}\cabecalho{S2\_02 -- Classes, Funções e Chamadas}{}
\bloco{O que o card pede} Modelar \codigo{Ativo} e \codigo{Vulnerabilidade} com classes e organizar as operações em funções reutilizáveis.
\bloco{Conceito} Uma classe é uma forma de guardar, no mesmo lugar, os dados que descrevem uma entidade e ações simples dela. Por exemplo, um objeto \codigo{Ativo} possui \codigo{id}, \codigo{nome}, \codigo{criticidade}, uma lista de vulnerabilidades e histórico. A função \codigo{para\_dicionario()} prepara esse objeto para o JSON. Uma função recebe dados, faz uma tarefa previsível e devolve um resultado.
\bloco{Aplicação} \codigo{modelos.py} declara as duas classes. \codigo{inventario.py} contém o CRUD sem \codigo{input} ou \codigo{print}; por isso a mesma regra pode ser chamada pelo menu ou pelo teste. Exemplo: \codigo{cadastrar\_ativo(ativos, 10, "SRV", "TI", "DC", 2, "Alta")} devolve o objeto criado ou informa um erro de regra.
\bloco{Por que separar} O terminal não decide se o CVSS é válido; ele apenas pergunta o valor. A regra \codigo{validar\_cvss} decide o intervalo. Isso evita duplicar a mesma lógica e deixa o resultado previsível.
\bloco{Checklist atendido} Classes do domínio; retornos definidos; ausência de globais desnecessárias; I/O fora das regras; exemplo coberto por teste sem erro.
\end{document}
